\documentclass[11pt]{article}

\usepackage[margin=1in]{geometry}
\usepackage{amsmath,amssymb,amsthm}
\usepackage{mathtools}
\usepackage{enumitem}
\usepackage{hyperref}

\theoremstyle{plain}
\newtheorem{theorem}{Theorem}
\newtheorem{lemma}[theorem]{Lemma}
\newtheorem{corollary}[theorem]{Corollary}
\newtheorem{definition}[theorem]{Definition}
\newtheorem{remark}[theorem]{Remark}

% ---------------------------------------------------------------------
% Macros
% ---------------------------------------------------------------------
\newcommand{\Ssem}[1]{\mathsf{S}[\![#1]\!]}
\newcommand{\Lsem}[1]{\mathsf{L}[\![#1]\!]}
\newcommand{\Vsem}[1]{\mathsf{V}[\![#1]\!]}

\newcommand{\Fm}[1]{F_{#1}}
\newcommand{\Rm}[1]{\widehat{R}_{#1}}
\newcommand{\Sm}[1]{\widehat{S}_{#1}}
\newcommand{\Stepm}[1]{\widehat{\mathit{Step}}_{#1}}

\newcommand{\eps}{\epsilon}
\newcommand{\op}{\mathit{op}}
\newcommand{\then}{\mathrel{\blacktriangleright}}
\newcommand{\loc}[1]{\langle #1 \rangle} % local construct, decorated separately
\newcommand{\tell}{\mathsf{tell}}
\newcommand{\censor}{\mathsf{censor}_0}
\newcommand{\up}{\mathsf{up}}
\newcommand{\glocalOp}{\mathsf{glocal}}
\newcommand{\rpt}[2]{\lfloor #1\rfloor_{#2}}
\newcommand{\lC}{\lambda_C}

\title{A Denotational Semantics for $\lC$ with Parameterized Handlers}
\author{}
\date{}

\begin{document}
\maketitle

\begin{abstract}
We give a denotational semantics for a fragment of the effect calculus $\lC$: an
expression language with algebraic effect operations, a distinguished \texttt{loss}
effect used to record real-valued costs, and effect handlers
$\mathtt{with}\ h\ \mathtt{from}\ e_1\ \mathtt{handle}\ e_2$ that thread a running
\emph{parameter} across successive resumptions of a handled computation's delimited
continuation. Types are interpreted by a family of selection functions
$\Sm\eps(X) = (X\to\Rm\eps)\to R\times\Stepm\eps(X)$, built over an ordinary, loss-free
free monad $\Fm\eps$ of ground-typed effect trees. Unlike a construction that attaches a
loss contribution to every node of a single tree monad, here the accumulated loss is
produced exactly once at each step of unfolding a selection function, entirely separately
from the pending-computation tree, and a node's own children are themselves selection
functions rather than bare trees. We develop $\Fm\eps$, $\Sm\eps$, the semantics of
expressions, and, in particular, the semantics of parameterized handlers, which we define
by a direct structural recursion on a selection function's own step-by-step unfolding --
dispatching on whether the next pending operation belongs to the handler being applied --
rather than by extending an algebra map homomorphically. We motivate this choice: once
handlers carry a running parameter, a clause may resume its own delimited continuation at
a parameter different from the one under which the handler was itself invoked, and a
nested handler's own choice continuation may leak a value computed under such a
resumption into an enclosing computation; together these defeat the uniform-in-the-
parameter homomorphism property that an extension-based account would need, whereas a
direct recursion on one pending step at a time requires no such property.
\end{abstract}

\section{Introduction}

$\lC$ interprets terms via a family of \emph{selection functions}
\[
  \Sm\eps(X) = (X\to \Rm\eps)\to R\times\Stepm\eps(X),
\]
for $R$ the (additive) monoid in which losses are recorded. Intuitively, a value of
$\Sm\eps(X)$ takes a \emph{loss continuation} $\gamma:X\to\Rm\eps$ -- an estimate of the
eventual loss incurred by each possible result -- and, run against it, produces a pair: the
loss $r\in R$ incurred at this one step of the computation, together with either a final
value or one further pending operation call, whose own continuations are again selection
functions run against the same $\gamma$. Two design choices distinguish this construction
from one in which the accumulated loss is instead attached to every node of a single
effect-tree monad, with handler semantics obtained as the (unique) homomorphic extension of
an algebra map:

\begin{enumerate}
\item The accumulated loss is kept entirely separate from the pending-computation tree, as
  a plain pair produced once at each step of unfolding a selection function, rather than
  woven into the tree's own recursive structure. The auxiliary monad $\Fm\eps$ used to
  represent finished loss \emph{reports} is consequently an entirely ordinary, loss-free
  free monad: no loss is attached to any of its own nodes, at all.
\item Handler semantics is defined by a direct structural recursion on a selection
  function's own step-by-step unfolding, dispatching on whether the next pending operation
  belongs to the handler being applied, rather than by first building an algebra structure
  and appealing to its universal (homomorphic-extension) property.
\end{enumerate}

The second choice matters once handlers carry a running parameter. A handler's clause may
resume its own delimited continuation at a parameter \emph{different} from the one under
which the handler was itself invoked -- choosing, at each resumption, both a fresh
parameter value and a point to resume at. Independently, a nested handler's own choice
continuation -- the value it reports back describing what loss it \emph{would} incur, were
a given resumption selected -- may itself leak a value computed under such a resumption
into an enclosing computation. Together, these two features defeat the
uniform-in-the-parameter algebra-homomorphism property that an extension-based account of
handler semantics would need: the map obtained by extending an algebra homomorphically is
only guaranteed to behave well when it is applied \emph{uniformly}, and a resumption at a
freshly chosen parameter is precisely a non-uniform use of the extension. A direct
recursion on one pending step at a time never needs to state, let alone prove, any such
uniform property in the first place: each case of the recursion only ever inspects the
selection function's immediately next step, dispatches or recurses into a strictly smaller
subtree, and the whole definition is well founded by ordinary structural recursion on that
unfolding -- mirroring, step for step, how a handled computation is itself reduced
operationally, one pending operation at a time.

As with any construction that restricts attention to \emph{ground} (first-order) operation
signatures -- no operation may take or return a value of function type, though the object
language's expressions remain fully higher-order throughout, and handler clauses may
themselves return functions -- this keeps the auxiliary monad $\Fm\eps$ an ordinary,
non-mutually-recursive inductive definition, needing only the ground-type interpretation
of types, not the general one (which is defined mutually with $\Sm\eps$ itself, via the
interpretation of function types).

Section~2 recalls the syntax and semantics of types, including parameterized handlers.
Section~3 develops the monads $\Fm\eps$ and $\Sm\eps$. Section~4 gives the semantics of
expressions. Section~5 gives the semantics of handlers.

\section{Syntax and Semantics of Types}

Types $\sigma,\tau$ are built from ground (first-order) base types, tuples, and function
types $\sigma\to\tau\,!\,\eps$ decorated with an effect context $\eps$ (a multiset of
effect labels available to the function body, written additively: $\eps(\ell)$ counts the
occurrences of $\ell$ in $\eps$, and $\eps\ell$ denotes $\eps$ extended by one further,
outermost occurrence of $\ell$). A signature $\Sigma$ fixes, for each effect label $\ell$, a
family of operations $\op:\mathit{out}\xrightarrow{\ell}\mathit{in}$ with
$\mathit{out},\mathit{in}$ ground types, and a distinguished ground type $\mathtt{loss}$ of
loss constants $r\in R$.

Values include variables, ground constants (base constants, loss constants $r:\mathtt{loss}$,
and $()$), pairs, and abstractions $\lambda^\eps x{:}\sigma.\,e$. Expressions $e$ include
values, applications of a primitive (first-order) function, tuples, projections, function
application, an operation call $\op(e)$ (typed against a witness that $\ell\in\eps$),
$\mathtt{loss}(e)$ to record a loss, a $(\then)$-construct $e_1\then\lambda^{\eps_1}
x{:}\sigma_1.\,e_2$ that sequences a hypothetical loss continuation, a local construct
$\loc{e}_g^{\eps_1}$ that substitutes a different loss continuation while evaluating $e$,
$\mathtt{reset}\ e$ to discard recorded loss, and handling
$\mathtt{with}\ h\ \mathtt{from}\ e_1\ \mathtt{handle}\ e_2$.

A handler $h$ for effect label $\ell$, parameter type $\mathit{par}$, body type $\sigma$,
and result type $\sigma'$ at ambient effect context $\eps$, supplies, for each operation
$\op:\mathit{out}\xrightarrow{\ell}\mathit{in}$, a \emph{clause} $e_\op$, typed in a context
extended by four variables -- the \emph{running parameter} $z_p{:}\mathit{par}$, the
operation's own argument $z_v{:}\mathit{out}$, a \emph{choice continuation}
$l:(\mathit{par}\times\mathit{in})\Rightarrow\mathtt{loss}\,!\,\eps$, and a \emph{delimited
continuation} $k:(\mathit{par}\times\mathit{in})\Rightarrow\sigma'\,!\,\eps$ -- of result
type $\sigma'\,!\,\eps$; and a \emph{return clause} $e_r$, typed in a context extended by
the running parameter $z_p{:}\mathit{par}$ and a result variable $z_x{:}\sigma$, likewise
of type $\sigma'\,!\,\eps$. Both $l$ and $k$ take a \emph{pair} $(p',a)$ of a freshly chosen
parameter value and an operation-result value: a clause resumes the handled computation by
supplying, at each resumption, both the parameter the handler should continue running
under and the point in the operation's result type to resume at. In
$\mathtt{with}\ h\ \mathtt{from}\ e_1\ \mathtt{handle}\ e_2$, the expression
$e_1:\mathit{par}\,!\,\eps$ supplies the \emph{initial} parameter value under which $h$
begins running $e_2:\sigma\,!\,\eps\ell$.

The semantics of types, $\Ssem\sigma$, is given by
\[
  \Ssem{b} = [\![b]\!], \qquad
  \Ssem{(\sigma_1,\dots,\sigma_n)} = \Ssem{\sigma_1}\times\cdots\times\Ssem{\sigma_n}, \qquad
  \Ssem{\sigma\to\tau\,!\,\eps} = \Ssem{\sigma}\to \Sm\eps(\Ssem{\tau}),
\]
where $\Sm\eps$ is the selection-function family developed in Section~3.

\section{The Monads $\Fm\eps$ and $\Sm\eps$}

\subsection{The carrier $\Fm\eps$}

An $\eps$-algebra is a set $X$ with maps
$\varphi_{\ell,\op,i} : \Ssem{\mathit{out}}\times X^{\Ssem{\mathit{in}}}\to X$ for each
$\ell\in\eps$, $\op:\mathit{out}\xrightarrow{\ell}\mathit{in}$, $0<i\le\eps(\ell)$; the free
such algebra on a set $X$ is written $\Fm\eps(X)$.

\begin{definition}[The monad $\Fm\eps$]
For a set $X$, let $\Fm\eps(X)$ be the least set $Y$ such that
\[
  Y \;=\; X \;+\; \sum_{\substack{\ell\in\eps,\ \op:\mathit{out}\xrightarrow{\ell}\mathit{in}\\ 0<i\le\eps(\ell)}}
        \Ssem{\mathit{out}} \times Y^{\Ssem{\mathit{in}}}.
\]
We write an element as either a leaf $\iota_1(x)$ for $x\in X$, or an internal node
$\iota_2\big((\ell,\op,i),(o,\kappa)\big)$; $\Fm\eps$ is the ordinary free monad for the
$\eps$-signature, with unit $\eta^{\Fm\eps}(x)=\iota_1(x)$, algebra structure
$\varphi_{\ell,\op,i}(o,\kappa)=\iota_2((\ell,\op,i),(o,\kappa))$, and Kleisli extension
$f\dg{}:\Fm\eps(X)\to\Fm\eps(Y)$ for $f:X\to\Fm\eps(Y)$ given, as usual, by structural
recursion:
\[
  f\dg{}(\iota_1(x)) = f(x), \qquad
  f\dg{}\big(\iota_2((\ell,\op,i),(o,\kappa))\big) = \iota_2\big((\ell,\op,i),\,(o,\ f\dg{}\circ\kappa)\big).
\]
No loss is attached to any node of $\Fm\eps(X)$ at all: unlike a construction in which the
accumulated loss is embedded throughout a single tree monad, here $\Fm\eps$ is kept
entirely free of $R$, and is used solely to represent \emph{finished} computations --
values, or trees of pending operation calls with no further loss bookkeeping of their own.
\end{definition}

For $\mathit{sub}:\eps_1\subseteq\eps$ (every occurrence available in $\eps_1$ is also
available in $\eps$), write $(\mathit{sub})_*:\Fm{\eps_1}(X)\to\Fm\eps(X)$ for the evident
structural relabelling of occurrence witnesses, leaving values and tree shape untouched.

We write $\Rm\eps := \Fm\eps(R)$ for the type of \emph{loss reports}: a finished
computation whose leaves carry an accumulated loss value rather than an ordinary result.

\subsection{Selection functions $\Sm\eps$}

\begin{definition}[The mutual family $\Sm\eps,\Stepm\eps$]
For a set $X$, let $\Sm\eps(X)$ and $\Stepm\eps(X)$ be given by the simultaneous least
solution of
\[
  \Sm\eps(X) \;=\; (X\to\Rm\eps)\to R\times\Stepm\eps(X), \qquad
  \Stepm\eps(X) \;=\; X \;+\; \sum_{\substack{\ell\in\eps,\ \op:\mathit{out}\xrightarrow{\ell}\mathit{in}\\ 0<i\le\eps(\ell)}}
        \Ssem{\mathit{out}} \times \Sm\eps(X)^{\Ssem{\mathit{in}}}.
\]
As with $\Fm\eps$ we write an element of $\Stepm\eps(X)$ as either $\mathrm{leaf}(x)$ or
$\mathrm{node}\big((\ell,\op,i),(o,\kappa)\big)$; since $\Sm\eps(X)$ is a function type, we
apply it directly, writing $p(\gamma)$ for a selection function $p$ run against a loss
continuation $\gamma$.
\end{definition}

\begin{remark}
$\Stepm\eps(X)$ is defined by exactly the same equation as one unfolding of $\Fm\eps$'s
own defining equation -- except that its recursive occurrences are of $\Sm\eps(X)$ itself,
not of $\Stepm\eps(X)$. A node's children are therefore \emph{further selection
functions}, not bare trees: unfolding one step of $p\in\Sm\eps(X)$ against $\gamma$ yields
either a finished value or one pending operation call, whose continuations must again be
run against a (possibly different) loss continuation to make further progress. This is the
structural feature that makes $\Sm\eps$ genuinely different from a family built as
$(X\to\Rm\eps)\to \Fm\eps(R\times X)$ or similar: there, once a loss continuation is
supplied, the resulting tree is an ordinary, loss-free value of a plain monad, independent
of $\Sm\eps$ itself. Here the recursive structure of $\Sm\eps$ and $\Stepm\eps$ never
bottoms out into a monad independent of $\Sm\eps$; the accumulated loss $R$ that
\emph{does} appear is produced once per step of the unfolding, sitting entirely outside
$\Stepm\eps(X)$, rather than being threaded through the recursive datatype at all.

One consequence: because a selection function's children are again selection functions
rather than plain trees, $\Sm{\eps_1}(X)$ cannot simply be reinterpreted as an instance of
$\Sm\eps(X)$ the way $\Fm{\eps_1}(X)$ can be relabelled into $\Fm\eps(X)$ by
$(\mathit{sub})_*$: changing the ambient effect context changes the domain of every loss
continuation nested inside it. Section~3.3 gives the genuine combinator, $\glocalOp$, that
this requires.
\end{remark}

\subsection{The combinators of $\Sm\eps$}

\paragraph{Unit and operation calls.}
\[
  \eta^{\Sm\eps}(x)(\gamma) = (0,\,\mathrm{leaf}(x)), \qquad
  \varphi_{\ell,\op,i}(o,\kappa)(\gamma) = (0,\,\mathrm{node}((\ell,\op,i),(o,\kappa))).
\]

\paragraph{Collapsing a selection function to a report.} For $p\in\Sm\eps(X)$ and a loss
continuation $\gamma:X\to\Rm\eps$, define $\rpt{p}{\gamma}\in\Rm\eps$, the result of running
$p$ against $\gamma$ \emph{throughout} -- at $p$ itself and at every one of its own pending
continuations -- and collapsing the result into a single finished report:
\[
  \rpt{p}{\gamma} =
  \begin{cases}
    \Fm\eps(r+{-})\big(\gamma(x)\big) & p(\gamma) = (r,\mathrm{leaf}(x)) \\[2pt]
    \Fm\eps(r+{-})\Big(\varphi_{\ell,\op,i}\big(o,\ \lambda a.\ \rpt{\kappa(a)}{\gamma}\big)\Big) & p(\gamma) = (r,\mathrm{node}((\ell,\op,i),(o,\kappa))),
  \end{cases}
\]
where $\Fm\eps(r+{-})$ denotes the functorial action of $\Fm\eps$ applying $r+{-}$ at every
leaf. That is: at a finished value, add this step's own loss $r$ into whatever $\gamma$
reports there; at a pending operation, leave the node itself untouched, add $r$ into every
leaf beneath it, and recurse into each child under the same $\gamma$.

\paragraph{Kleisli extension.} For $f:X\to\Sm\eps(Y)$, define
$f\dg{}:\Sm\eps(X)\to\Sm\eps(Y)$, writing $\mathsf{let}_{\Sm\eps}\ x\in X\ \mathsf{be}\ p\
\mathsf{in}\ f(x) := f\dg{}(p)$, by first running
\[
  p\big(\lambda x.\,\rpt{f(x)}{\gamma}\big) = (r_1,w)
\]
and then, by cases on $w$,
\[
  \big(\mathsf{let}_{\Sm\eps}\ x\in X\ \mathsf{be}\ p\ \mathsf{in}\ f(x)\big)(\gamma) = (r_1+r_2,\,s_2), \qquad (r_2,s_2)=f(x')(\gamma), \qquad \text{if } w=\mathrm{leaf}(x'),
\]
\begin{multline*}
  \big(\mathsf{let}_{\Sm\eps}\ x\in X\ \mathsf{be}\ p\ \mathsf{in}\ f(x)\big)(\gamma) = \big(r_1,\ \mathrm{node}((\ell,\op,i),(o,\lambda a.\ \mathsf{let}_{\Sm\eps}\ x\in X\ \mathsf{be}\ \kappa(a)\ \mathsf{in}\ f(x)))\big), \\
  \text{if } w=\mathrm{node}((\ell,\op,i),(o,\kappa)).
\end{multline*}
That is: run $p$ under the hypothetical continuation $\lambda x.\rpt{f(x)}{\gamma}$ (what
$f$ would eventually report, at each possible $x$, run to completion against $\gamma$); if
$p$ is already finished at some $x'$, the whole computation continues as $f(x')$ run
directly against $\gamma$, with $p$'s own step-loss $r_1$ added on; if $p$ has a pending
operation, that operation is left exactly in place -- with loss $r_1$ untouched -- and only
its children are rebound.

\paragraph{Adding loss.} For $r\in R$, $\tell(r):\Sm\eps(X)\to\Sm\eps(X)$ adds $r$ into a
selection function's own next step, wherever that step occurs:
\[
  \tell(r)(p)(\gamma) = \big(r+r_1,\ s_1\big), \qquad (r_1,s_1) = p(\gamma).
\]

\paragraph{Promoting a report.} $\up:\Fm\eps(X)\to\Sm\eps(X)$ promotes a finished,
$\gamma$-independent report into a selection function that ignores whatever continuation it
is later run against:
\[
  \up(\iota_1(x))(\gamma) = (0,\mathrm{leaf}(x)), \qquad
  \up\big(\iota_2((\ell,\op,i),(o,\kappa))\big)(\gamma) = \big(0,\ \mathrm{node}((\ell,\op,i),(o,\lambda a.\,\up(\kappa(a))))\big).
\]

\paragraph{Fixing a loss continuation.} For $\mathit{sub}:\eps_1\subseteq\eps$ and
$\gamma_1:X\to\Rm{\eps_1}$, define $\glocalOp(\mathit{sub},\gamma_1,-):\Sm{\eps_1}(X)\to\Sm\eps(X)$,
which runs $p$ against the fixed continuation $\gamma_1$ throughout -- ignoring whatever
$\gamma$ it is itself later run against -- while relabelling every occurrence witness along
the way via $\mathit{sub}$:
\[
  \glocalOp(\mathit{sub},\gamma_1,p)(\gamma) = (r,\mathrm{leaf}(x)), \qquad \text{if } p(\gamma_1) = (r,\mathrm{leaf}(x)),
\]
\begin{multline*}
  \glocalOp(\mathit{sub},\gamma_1,p)(\gamma) = \big(r,\ \mathrm{node}(\mathit{sub}(\ell,\op,i),(o,\lambda a.\ \glocalOp(\mathit{sub},\gamma_1,\kappa(a))))\big), \\
  \text{if } p(\gamma_1) = (r,\mathrm{node}((\ell,\op,i),(o,\kappa))).
\end{multline*}

\paragraph{Discarding recorded loss.} $\censor:\Sm\eps(X)\to\Sm\eps(X)$ recursively zeroes
the loss recorded at every step of $p$'s own unfolding (run against a caller-supplied
$\gamma$), leaving values and tree shape untouched:
\[
  \censor(p)(\gamma) = (0,\mathrm{leaf}(x)), \qquad \text{if } p(\gamma) = (r,\mathrm{leaf}(x)),
\]
\begin{multline*}
  \censor(p)(\gamma) = \big(0,\ \mathrm{node}((\ell,\op,i),(o,\lambda a.\,\censor(\kappa(a))))\big), \\
  \text{if } p(\gamma) = (r,\mathrm{node}((\ell,\op,i),(o,\kappa))).
\end{multline*}

\section{Semantics of Expressions}

\subsection{Values}

\[
  \Vsem{x}(\rho) = \rho(x), \qquad
  \Vsem{c}(\rho) = [\![c]\!], \qquad
  \Vsem{(v,w)}(\rho) = (\Vsem{v}(\rho),\Vsem{w}(\rho)),
\]
\[
  \Vsem{\lambda^\eps x{:}\sigma.\,e}(\rho) = \lambda a.\ \Ssem{e}(\rho[a/x]).
\]

\subsection{The auxiliary loss-function semantics $\mathcal L$}

For $\Gamma,x{:}\sigma\vdash e:\mathtt{loss}\,!\,\eps$, define
$\Lsem{\lambda^\eps x{:}\sigma.\,e}:\Sm\Gamma\to\Sm\sigma\to\Rm\eps$ by
\[
  \Lsem{\lambda^\eps x{:}\sigma.\,e}(\rho)(a) = \rpt{\Vsem{\lambda^\eps x{:}\sigma.\,e}(\rho)(a)}{\eta^{\Fm\eps}}.
\]

\begin{remark}
The loss function $g=\lambda^\eps x{:}\sigma.\,e$ is itself an ordinary value of function
type $\sigma\to\mathtt{loss}\,!\,\eps$; its own body may call operations, sequence further
loss continuations, and eventually produce a value of type $\mathtt{loss}$, i.e.\ an $r\in
R$. That eventual value \emph{is itself} a further loss contribution -- $g$ is being run
purely to find out what loss it would report, so whatever it eventually returns must be
folded into the total, not discarded. Collapsing $\Vsem g(\rho)(a)$ under the loss
continuation $\eta^{\Fm\eps}:R\to\Rm\eps$ (which embeds a returned loss value as itself)
does exactly this, via $\rpt{-}{-}$'s own leaf case ($\Fm\eps(r+{-})(\gamma(x))$, here with
$\gamma(x)=\eta^{\Fm\eps}(x)$). Collapsing instead under a continuation that
\emph{discards} its argument (e.g.\ $\lambda r.\,\eta^{\Fm\eps}(0)$) would silently drop
$g$'s own returned value whenever $g$'s body ends in something other than a bare
$\mathtt{loss}(-)$ call -- for instance a loss function that simply evaluates to a constant
$r_0:\mathtt{loss}$ without ever calling $\mathtt{loss}(-)$ explicitly would then be
indistinguishable from one reporting no loss at all, which is not the intended reading of
$\mathcal L$.
\end{remark}

\subsection{Core constructs}

\[
  \Ssem{v}(\rho) = \eta^{\Sm\eps}(\Vsem{v}(\rho)), \qquad
  \Ssem{f(e)}(\rho) = \mathsf{let}_{\Sm\eps}\ a\in\Ssem\gamma\ \mathsf{be}\ \Ssem e(\rho)\ \mathsf{in}\ \eta^{\Sm\eps}([\![f]\!](a)),
\]
\[
  \Ssem{(e_1,e_2)}(\rho) = \mathsf{let}_{\Sm\eps}\ a\in\Ssem{\sigma_1}\ \mathsf{be}\ \Ssem{e_1}(\rho)\ \mathsf{in}\ \mathsf{let}_{\Sm\eps}\ b\in\Ssem{\sigma_2}\ \mathsf{be}\ \Ssem{e_2}(\rho)\ \mathsf{in}\ \eta^{\Sm\eps}((a,b)),
\]
with $\mathtt{fst}$, $\mathtt{snd}$, and application $e_1\,e_2$ the evident generic clauses
of a monadic semantics ($\Ssem{v\,e}(\rho) = \mathsf{let}_{\Sm\eps}\ a\in\Ssem\sigma\
\mathsf{be}\ \Ssem e(\rho)\ \mathsf{in}\ \Vsem v(\rho)(a)$), and
\[
  \Ssem{\op(e)}(\rho) = \mathsf{let}_{\Sm\eps}\ a\in\Ssem{\mathit{out}}\ \mathsf{be}\ \Ssem{e}(\rho)
     \ \mathsf{in}\ \varphi_{\ell,\op,\eps(\ell)}\big(a,\,(\eta^{\Sm\eps})^{\Ssem{\mathit{in}}}\big).
\]

\paragraph{\texttt{loss(e)}.} A bare $\tell$ is directly available on $\Sm\eps$, so no
pattern matching on the shape of $\Ssem e$ is needed:
\[
  \Ssem{\mathtt{loss}(e)}(\rho) = \mathsf{let}_{\Sm\eps}\ r\in R\ \mathsf{be}\ \Ssem{e}(\rho)\ \mathsf{in}\ \tell(r)\big(\eta^{\Sm\eps}(())\big).
\]

\paragraph{$e_1\then\lambda^{\eps_1} x{:}\sigma_1.\,e_2$.} Collapse $e_1$ to a report using
$g=\lambda x.e_2$'s (widened) semantics as the loss continuation throughout, then promote
that report back into a selection function via $\up$:
\[
  \Ssem{e_1\then\lambda^{\eps_1}x{:}\sigma_1.\,e_2}(\rho) = \up\Big(\rpt{\Ssem{e_1}(\rho)}{\lambda a.\ (\mathit{sub})_*\big(\Lsem{\lambda x.\,e_2}(\rho)(a)\big)}\Big),
\]
for $\mathit{sub}:\eps_1\subseteq\eps$ (left implicit in the surface notation, as
throughout). Because $\up$ ignores whatever continuation it is itself later run against,
$\Ssem{e_1\then g}(\rho)$ does not depend on the ambient loss continuation the surrounding
computation eventually supplies -- $\then$ has already fixed its own loss accounting via
$g$ by the time it produces a value.

\paragraph{$\loc{e}_g^{\eps_1}$.} For $e:\sigma\,!\,\eps_1$ and $g:\sigma\to\mathtt{loss}\,!\,\eps_2$
with $\eps_2\subseteq\eps_1\subseteq\eps$ (again left implicit), this substitutes a
different loss continuation while evaluating $e$, using $\glocalOp$ (Section~3.3) rather
than a bare application, since $\Sm{\eps_1}(\sigma)$ is not literally reusable as a value of
$\Sm\eps(\sigma)$:
\[
  \Ssem{\loc{e}_g^{\eps_1}}(\rho) = \glocalOp\Big(\mathit{sub}_2,\ \lambda a.\ (\mathit{sub}_1)_*\big(\Lsem g(\rho)(a)\big),\ \Ssem e(\rho)\Big).
\]

\paragraph{\texttt{reset}\,$e$.} Discard all loss recorded while evaluating $e$, keeping
its value, via $\censor$:
\[
  \Ssem{\mathtt{reset}\ e}(\rho) = \censor\big(\Ssem e(\rho)\big).
\]

\paragraph{Handling.} Evaluate $e_1$ to obtain the initial parameter value, then dispatch to
the handler semantics developed in Section~5:
\[
  \Ssem{\mathtt{with}\ h\ \mathtt{from}\ e_1\ \mathtt{handle}\ e_2}(\rho) = \mathsf{let}_{\Sm\eps}\ p\in\Ssem{\mathit{par}}\ \mathsf{be}\ \Ssem{e_1}(\rho)\ \mathsf{in}\ \mathcal H_h(\rho,p,\Ssem{e_2}(\rho)).
\]

\section{Semantics of Handlers}

Fix a handler $h$ for effect label $\ell$, and an environment $\rho\in\Sm\Gamma$. We define
$\mathcal H_h(\rho,p,G) \in \Sm\eps(\Ssem{\sigma'})$ for a running parameter value
$p\in\Ssem{\mathit{par}}$ and a computation $G\in\Sm{\eps\ell}(\Ssem\sigma)$ still to be
handled, by direct structural recursion on $G$'s own next unfolding step -- \emph{not} by
building an $\eps\ell$-algebra structure and taking a homomorphic extension.

\paragraph{The return clause.} Write
\[
  s(\rho,p,a) = \Ssem{e_r}\big(\rho[p/z_p,a/z_x]\big) \in \Sm\eps(\Ssem{\sigma'})
\]
for the semantics of $h$'s return clause, at running parameter $p$ and result value $a$.

\paragraph{The handler semantics.} Running $\mathcal H_h(\rho,p,G)$ against a loss
continuation $\gamma:\Ssem{\sigma'}\to\Rm\eps$ first runs $G$ itself against the
hypothetical continuation $\lambda a.\ \iota\big(\rpt{s(\rho,p,a)}{\gamma}\big)$ -- what the
return clause would eventually report, at each possible result $a$, were $G$ to finish
there -- where $\iota:\Rm\eps\hookrightarrow\Rm{\eps\ell}$ is the evident inclusion induced
by $\eps\subseteq\eps\ell$ (needed since $G$ itself lives at the wider context $\eps\ell$).
The resulting step is then dispatched by $\Delta_h(\rho,\gamma,p)$:
\[
  \mathcal H_h(\rho,p,G)(\gamma) = \Delta_h(\rho,\gamma,p)\Big(G\big(\lambda a.\ \iota(\rpt{s(\rho,p,a)}{\gamma})\big)\Big).
\]

\paragraph{Dispatch.} $\Delta_h(\rho,\gamma,p)$ takes a step $(r,w)\in R\times
\Stepm{\eps\ell}(\Ssem\sigma)$ of $G$'s own unfolding to a step of $\Sm\eps(\Ssem{\sigma'})$,
by cases on $w$:
\begin{itemize}
\item \emph{$G$ is already finished}, $w=\mathrm{leaf}(x)$: invoke the return clause at the
  running parameter and the produced value, and combine the loss $r_2$ this incurs with
  whatever loss $r$ $G$ had already accumulated on the way to $x$:
  \[
    \Delta_h(\rho,\gamma,p)(r,\mathrm{leaf}(x)) = (r+r_2,\,s_2), \qquad (r_2,s_2) = s(\rho,p,x)(\gamma).
  \]
\item \emph{$G$'s next step is an operation not handled by $h$} (i.e.\ at an occurrence
  already present in $\eps$, shallower than the newly introduced occurrence of $\ell$ that
  $G$ was typed against): pass the step through unchanged, demoting its occurrence witness,
  and recurse into each of its continuations:
  \begin{multline*}
    \Delta_h(\rho,\gamma,p)\big(r,\mathrm{node}((\ell_1,\op,i),(o,\kappa))\big) \\
      = \big(r,\ \mathrm{node}((\ell_1,\op,i),(o,\lambda a.\ \mathcal H_h(\rho,p,\kappa(a))))\big), \qquad \ell_1\ne \ell.
  \end{multline*}
\item \emph{$G$'s next step is one of $h$'s own operations} $\op_j$, at the newly
  introduced occurrence: dispatch directly to the corresponding object-language clause,
  instantiated at the running parameter, the operation's own argument, and the choice and
  delimited continuations $l_1,k_1$ built from $\kappa$ below, and run it against $\gamma$;
  again combine its own freshly incurred loss $r_2$ with $G$'s own accumulated $r$:
  \begin{multline*}
    \Delta_h(\rho,\gamma,p)\big(r,\mathrm{node}(\op_j,(o,\kappa))\big) = (r+r_2,\,s_2), \qquad \\
    (r_2,s_2) = \Ssem{e_j}\big(\rho[p/z_p,\,o/z_v,\,l_1/l,\,k_1/k]\big)(\gamma).
  \end{multline*}
\end{itemize}

\paragraph{Choice and delimited continuations.} At each freshly chosen parameter $p'$ and
resumption point $a$ offered by $\kappa$, the delimited continuation $k_1$ resumes the same
handler, running under the caller's own ambient loss continuation $\gamma$ throughout (via
$\glocalOp$ at the identity inclusion); the choice continuation $l_1$ reports what loss that
same resumption \emph{would} incur, collapsed to a bare report and promoted back into a
selection function that ignores whatever continuation it is itself later run against (via
$\up\circ\rpt{-}{\gamma}$), so that it can be used as an ordinary value of type
$\mathtt{loss}$ inside the clause body:
\[
  k_1(p',a) = \glocalOp\big(\mathrm{id},\ \gamma,\ \mathcal H_h(\rho,p',\kappa(a))\big), \qquad
  l_1(p',a) = \up\Big(\rpt{\mathcal H_h(\rho,p',\kappa(a))}{\gamma}\Big).
\]

\begin{remark}
Every recursive call above -- inside $\Delta_h$'s pass-through case, and inside $k_1,l_1$ --
is made on $\kappa(a)$, one of $G$'s own children, for some $a$ chosen by $G$ itself: the
recursion is on $G$'s own unfolding, one pending step at a time, and never on the running
parameter $p$. This is exactly what lets $k_1$ and $l_1$ resume the handler at a freshly
chosen $p'$, different at each call, with no obligation to relate the behaviour of
$\mathcal H_h(\rho,p,-)$ and $\mathcal H_h(\rho,p',-)$ as $p$ ranges over
$\Ssem{\mathit{par}}$: nothing in the definition is stated, or needs to be proved, uniformly
in $p$. An account built instead by first fixing a single algebra structure on
$\Rm\eps(\Ssem{\sigma'})^{\Ssem{\mathit{par}}}$ and then extending a generator map
homomorphically would need exactly such a uniform property -- the extension's own defining
property is a statement about the map's behaviour at \emph{every} parameter value
simultaneously -- which a clause resuming at a freshly chosen $p'\ne p$, combined with a
nested handler's own choice continuation reporting a value computed under such a
resumption, need not respect.
\end{remark}

\end{document}
